\section{The finite initial affine certificate}
\label{app:finite-seed}

This appendix records the only finite exceptional reciprocal calculation used
by the affine-ray provider.  It is proof data rather than part of the
categorical structure.

Put
\[
P(m)=\frac1m+\frac1{m+1}.
\]
At the initial cyclic stage we use the empty support for residue $0$ and the
following five binary states:
\[
\frac16=
\sum_{m\in\{35,56,90,119,152,170,455,494,714\}}P(m),
\tag{S1}
\]
\[
\frac13=
\sum_{m\in\{17,35,44,84,104,119,132,170,272,285,455,494,560,714,935\}}P(m),
\tag{S2}
\]
\[
\frac12=
\sum_{m\in\{9,17,34,44,77,84,90\}}P(m),
\tag{S3}
\]
\[
\frac23=
\sum_{m\in\{7,16,34,44,51,65,77,119,132,153,170,272,285,455,494\}}P(m),
\tag{S4}
\]
\[
\frac56=P(2).
\tag{S5}
\]
The corresponding component counts are
\[
0,9,15,7,15,1.
\tag{S6}
\]
Every displayed state is a union of pairwise non-adjacent binary blocks.  In
each start set the difference between consecutive starts is at least $3$.

The identities (S1)--(S5) are finite rational equalities.  For example, direct
reduction gives
\[
\frac{19399380}{116396280}=\frac16,
\qquad
\frac{77597520}{232792560}=\frac13,
\qquad
\frac{1531530}{3063060}=\frac12,
\]
\[
\frac{155195040}{232792560}=\frac23,
\qquad
\frac56=\frac56.
\]
Thus the six states represent the complete affine residue fibre
\[
\left\{0,\frac16,\frac13,\frac12,\frac23,\frac56\right\}
\subset\Q/\Z
\]
at the finite root stage.  Their maximum component count is a finite integer;
denote it by $M_*$.  Its numerical value is never used above the finite
certificate.

The finitely many bridge rows which are not yet covered by the asymptotic
strictness argument have lower bases
\[
12,20,30,60,140,210,420,840.
\]
For these rows the nearest relevant seed starts and minimum support-vertex
distances are
\[
\begin{array}{c|cccccccc}
a&12&20&30&60&140&210&420&840\\ \hline
m_{\rm near}&9&17&34&56&132&170&455&935\\
\min |u-v|&2&2&2&3&7&39&33&93.
\end{array}
\]
Hence no cross edge occurs in the reflexive successor graph.  All subsequent
bridge rows are covered by the uniform separation estimate of
Lemma~\ref{lem:complete-bridge}.

Finally we verify that the finite seed leaves an open reciprocal-value margin.
Every seed branch has mass at most $5/6$.  A later complete bridge with lower
modulus $D$ has mass less than $6/D$.  Beyond the initial stage the first lower
modulus is $60$, and each subsequent nontrivial compact jump multiplies the
lower modulus by a prime at least $2$.  Therefore the complete future bridge
tail has mass less than
\[
\frac6{60}\sum_{j\ge0}2^{-j}=\frac15.
\]
Every seed-plus-future-bridge branch consequently satisfies
\[
W<\frac56+\frac15=rac{31}{30}<2.
\]
Thus a positive open literalization margin remains even before the later
neutral equalization coordinates are chosen.  Since those coordinates can be
made arbitrarily light, the Prefix construction of
Section~\ref{sec:affine-provider} follows.
